\icode{E-MI-SI-L2-261107-002}
实验小组为探究弹簧弹力与形变量之间关系，选取了 \(a\)，\(b\) 两条弹簧，通过实验操作、数据记录和处理，描绘出弹力与弹簧长度的关系如图所示，下列说法正确的是\paren[]
	
\begin{tikzpicture}[>=stealth,font=\scriptsize]
	\draw[->,blue] (0,0)node[below left]{ \(O\)}--(3,0)node[right]{ \(l\)};
	\draw[->,blue] (0,0)--(0,2.5)node[right]{ \(F\)};
	\draw[magenta] (0.6,0)node[below,blue]{ \(l_1\)}--++(70:2)node[left]{ \(a\)} (1,0)node[below,blue]{ \(l_2\)}--++(50:2)node[right]{ \(b\)};
\end{tikzpicture}
\begin{choices}
	\item 弹簧 \(a\) 的原长大于弹簧 \(b\)
	\item 弹簧 \(a\) 的劲度系数小于弹簧 \(b\)
	\item 在弹性限度内，用大小相等的力 \(F\) 拉伸弹簧，弹簧 \(b\) 的伸长量大于弹簧 \(a\)
	\item 在弹性限度内，用大小相等的力 \(F\) 压缩弹簧，弹簧 \(b\) 的压缩量小于弹簧 \(a\)
\end{choices}
\begin{answer}
    C\qquad 提示：
\end{answer}